%% kapa: exam=1 round=22 year=2011 period=1 form=A subject=economics subject_ko=경제원론 seq=28 number=68
%% source: Q-Net 감정평가사 기출문제 (cst00309) artlSeq=5036790 file=extracted/1008/5036790/2011년 제22회 감정평가사 1차시험 기출문제/1교시.hwp
%% conversion: kordoc 4.24.51 -> tools/gen1.py (auto); review=proofread-sonnet; flags=none
\begin{question}{68}
다음 중 옳은 것만을 모두 고른 것은?
\begin{framedtext}
ㄱ. 프리드만(M. Friedman)은 통화량을 일정률로 증가시키는 통화준칙을 주장한다. \par
ㄴ. 새고전학파(New Classical School)는 예측되는 정책은 항상 긍정적인 효과가 있다고 주장한다. \par
ㄷ. 새케인즈학파(New Keynesian School) 이론 중에는 메뉴비용(menu cost)의 존재로 총수요관리정책이 효과가 있다는 주장이 있다. \par
ㄹ. 실물경기변동론자들은 기술충격에 의한 총공급의 변동으로 경기변동을 설명한다. \par
ㅁ. 케인즈학파(Keynesian School)는 총공급의 변동이 경기변동의 가장 중요한 원인이라고 주장한다.
\end{framedtext}
\begin{choices}
  \item ㄱ, ㄴ, ㅁ
  \item ㄱ, ㄴ, ㄷ
  \item ㄱ, ㄷ, ㄹ
  \item ㄴ, ㄷ, ㅁ
  \item ㄴ, ㄷ, ㄹ
\end{choices}
\end{question}
\answer{3}
